\begin{theorem}[The causal arrow has no return]\label{thm:causal-arrow-no-return}\lean{Obsmetrologia.CausalSets.Causal.ArrowOfTime.future_not_past}\leanok\uses{thm:causal-settling-transitive-closure}
The settled future relation is directed reachability. If $a$ is in the future of
$b$, then $b$ is not in the future of $a$: the past step is excluded by the
acyclicity of causal settling. This is the first arrow-of-time theorem. It is a
causal no-return statement, not an entropy theorem and not a continuum-time
reconstruction.
\end{theorem}

\begin{theorem}[Ordinal clocks require successor advance]\label{thm:ordinal-arrow-successor-advance}\lean{Obsmetrologia.CausalSets.Causal.ArrowOfTime.causalOrdinalArrow}\leanok\uses{thm:causal-arrow-no-return}
An ordinal-staged arrow of time packages a temporal relation, an `Ord`-valued
stage map, a proof that each temporal step advances at least to the successor
stage, and a no-return law. A causal set carries such an ordinal arrow only when
the stage map and successor-advance proof are supplied; the no-return component
then comes from causal settling itself.
\end{theorem}

\begin{theorem}[Ordinal-stage well-founded descent]\label{thm:ordinal-stage-well-founded-descent}\notready\uses{thm:ordinal-arrow-successor-advance}
An ordinal-staged temporal relation whose every step advances by successor
stage should support a well-founded descent principle for reasoning backward
along time. The carrier exists, but the project still needs the strict
well-foundedness theorem for `Ord` successor advance. Until that ordinal descent
theorem is formalized, this remains a red node rather than a Nat-ranked
surrogate.
\end{theorem}
